% Options for packages loaded elsewhere
% Options for packages loaded elsewhere
\PassOptionsToPackage{unicode}{hyperref}
\PassOptionsToPackage{hyphens}{url}
\PassOptionsToPackage{dvipsnames,svgnames,x11names}{xcolor}
%
\documentclass[
  letterpaper,
  DIV=11,
  numbers=noendperiod]{scrartcl}
\usepackage{xcolor}
\usepackage{amsmath,amssymb}
\setcounter{secnumdepth}{5}
\usepackage{iftex}
\ifPDFTeX
  \usepackage[T1]{fontenc}
  \usepackage[utf8]{inputenc}
  \usepackage{textcomp} % provide euro and other symbols
\else % if luatex or xetex
  \usepackage{unicode-math} % this also loads fontspec
  \defaultfontfeatures{Scale=MatchLowercase}
  \defaultfontfeatures[\rmfamily]{Ligatures=TeX,Scale=1}
\fi
\usepackage{lmodern}
\ifPDFTeX\else
  % xetex/luatex font selection
\fi
% Use upquote if available, for straight quotes in verbatim environments
\IfFileExists{upquote.sty}{\usepackage{upquote}}{}
\IfFileExists{microtype.sty}{% use microtype if available
  \usepackage[]{microtype}
  \UseMicrotypeSet[protrusion]{basicmath} % disable protrusion for tt fonts
}{}
\makeatletter
\@ifundefined{KOMAClassName}{% if non-KOMA class
  \IfFileExists{parskip.sty}{%
    \usepackage{parskip}
  }{% else
    \setlength{\parindent}{0pt}
    \setlength{\parskip}{6pt plus 2pt minus 1pt}}
}{% if KOMA class
  \KOMAoptions{parskip=half}}
\makeatother
% Make \paragraph and \subparagraph free-standing
\makeatletter
\ifx\paragraph\undefined\else
  \let\oldparagraph\paragraph
  \renewcommand{\paragraph}{
    \@ifstar
      \xxxParagraphStar
      \xxxParagraphNoStar
  }
  \newcommand{\xxxParagraphStar}[1]{\oldparagraph*{#1}\mbox{}}
  \newcommand{\xxxParagraphNoStar}[1]{\oldparagraph{#1}\mbox{}}
\fi
\ifx\subparagraph\undefined\else
  \let\oldsubparagraph\subparagraph
  \renewcommand{\subparagraph}{
    \@ifstar
      \xxxSubParagraphStar
      \xxxSubParagraphNoStar
  }
  \newcommand{\xxxSubParagraphStar}[1]{\oldsubparagraph*{#1}\mbox{}}
  \newcommand{\xxxSubParagraphNoStar}[1]{\oldsubparagraph{#1}\mbox{}}
\fi
\makeatother


\usepackage{longtable,booktabs,array}
\usepackage{calc} % for calculating minipage widths
% Correct order of tables after \paragraph or \subparagraph
\usepackage{etoolbox}
\makeatletter
\patchcmd\longtable{\par}{\if@noskipsec\mbox{}\fi\par}{}{}
\makeatother
% Allow footnotes in longtable head/foot
\IfFileExists{footnotehyper.sty}{\usepackage{footnotehyper}}{\usepackage{footnote}}
\makesavenoteenv{longtable}
\usepackage{graphicx}
\makeatletter
\newsavebox\pandoc@box
\newcommand*\pandocbounded[1]{% scales image to fit in text height/width
  \sbox\pandoc@box{#1}%
  \Gscale@div\@tempa{\textheight}{\dimexpr\ht\pandoc@box+\dp\pandoc@box\relax}%
  \Gscale@div\@tempb{\linewidth}{\wd\pandoc@box}%
  \ifdim\@tempb\p@<\@tempa\p@\let\@tempa\@tempb\fi% select the smaller of both
  \ifdim\@tempa\p@<\p@\scalebox{\@tempa}{\usebox\pandoc@box}%
  \else\usebox{\pandoc@box}%
  \fi%
}
% Set default figure placement to htbp
\def\fps@figure{htbp}
\makeatother





\setlength{\emergencystretch}{3em} % prevent overfull lines

\providecommand{\tightlist}{%
  \setlength{\itemsep}{0pt}\setlength{\parskip}{0pt}}



 


\renewcommand{\thesection}{S\arabic{section}}
\renewcommand{\thesubsection}{S\arabic{section}.\arabic{subsection}}
\usepackage{tabularx}
\newcommand{\trow}[3]{%
  \texttt{\allowbreak #1} & #2 & #3 \\
}
\newcommand{\trowfour}[4]{%
  \texttt{\allowbreak #1} & #2 & #3 & #4 \\
}
\KOMAoption{captions}{tableheading}
\makeatletter
\@ifpackageloaded{caption}{}{\usepackage{caption}}
\AtBeginDocument{%
\ifdefined\contentsname
  \renewcommand*\contentsname{Table of contents}
\else
  \newcommand\contentsname{Table of contents}
\fi
\ifdefined\listfigurename
  \renewcommand*\listfigurename{List of Figures}
\else
  \newcommand\listfigurename{List of Figures}
\fi
\ifdefined\listtablename
  \renewcommand*\listtablename{List of Tables}
\else
  \newcommand\listtablename{List of Tables}
\fi
\ifdefined\figurename
  \renewcommand*\figurename{Figure}
\else
  \newcommand\figurename{Figure}
\fi
\ifdefined\tablename
  \renewcommand*\tablename{Table}
\else
  \newcommand\tablename{Table}
\fi
}
\@ifpackageloaded{float}{}{\usepackage{float}}
\floatstyle{ruled}
\@ifundefined{c@chapter}{\newfloat{codelisting}{h}{lop}}{\newfloat{codelisting}{h}{lop}[chapter]}
\floatname{codelisting}{Listing}
\newcommand*\listoflistings{\listof{codelisting}{List of Listings}}
\makeatother
\makeatletter
\makeatother
\makeatletter
\@ifpackageloaded{caption}{}{\usepackage{caption}}
\@ifpackageloaded{subcaption}{}{\usepackage{subcaption}}
\makeatother
\usepackage{bookmark}
\IfFileExists{xurl.sty}{\usepackage{xurl}}{} % add URL line breaks if available
\urlstyle{same}
\hypersetup{
  pdftitle={Supplementary Material: Implementation Details of the Agent-Based Psychological Resilience Simulation},
  pdfauthor={Aly Lamuri},
  colorlinks=true,
  linkcolor={blue},
  filecolor={Maroon},
  citecolor={Blue},
  urlcolor={Blue},
  pdfcreator={LaTeX via pandoc}}


\title{Supplementary Material: Implementation Details of the Agent-Based
Psychological Resilience Simulation}
\author{Aly Lamuri}
\date{}
\begin{document}
\maketitle


\section{Phase Function Protocol}\label{phase-function-protocol}

\subsection{Purpose}\label{purpose}

Define the type contract that all phase functions must satisfy, enabling
modular composition and independent testing.

\begin{itemize}
\tightlist
\item
  \textbf{Frequency:} N/A (interface definition)
\item
  \textbf{Inputs:} N/A (type definitions)
\item
  \textbf{Outputs:} N/A (type definitions)
\end{itemize}

\subsection{Definitions}\label{definitions}

\textbf{PhaseFunction Protocol:}

\begin{verbatim}
FUNCTION phase(state, config, rng) → PhaseOutput
\end{verbatim}

Every phase function accepts (state, config, rng) and returns
PhaseOutput. This uniform signature enables sequential composition.

\textbf{PhaseOutput:}

\begin{verbatim}
PhaseOutput = {
    state_delta: {key: value}   // updates to apply to agent state
    observation: {key: value}   // logging data, no state effect
}
\end{verbatim}

state\_delta contains updates to apply via apply\_delta(). observation
contains metrics that do not affect simulation state.

\textbf{PhaseFrequency:}

\begin{verbatim}
PhaseFrequency = event_driven | daily
\end{verbatim}

event\_driven: runs per stress event or interaction (multiple times/day)
daily: runs once per day during consolidation loop

Reference:
\href{https://github.com/namakala/abm-simulation/blob/7e40a44f82da76b18910b774cd882c938bc79992/src/python/phases/interfaces.py\#L1-L85}{src/python/phases/interfaces.py:L1-L85}

\section{AgentState Schema}\label{agentstate-schema}

\subsection{Purpose}\label{purpose-1}

Single source of truth for all mutable agent variables. Passed between
phases; phases read from and write to this schema via
\texttt{state\_delta}.

\begin{itemize}
\tightlist
\item
  \textbf{Frequency:} N/A (data structure)
\item
  \textbf{Inputs:} N/A
\item
  \textbf{Outputs:} N/A
\end{itemize}

\subsection{Field Definitions}\label{field-definitions}

\textbf{Core state:}

\begin{tabularx}{\textwidth}{p{4cm}p{1.8cm}p{2cm}X}
\toprule
\textbf{Field} & \textbf{Type} & \textbf{Range} & \textbf{Description} \\
\midrule
\trowfour{baseline\_resilience}{float}{[0, 1]}{Initial resilience (fixed)}
\trowfour{resilience}{float}{[0, 1]}{Current resilience}
\trowfour{baseline\_affect}{float}{[-1, 1]}{Initial affect (fixed)}
\trowfour{affect}{float}{[-1, 1]}{Current affect}
\trowfour{resources}{float}{[0, 1]}{Psychological/physical resources}
\bottomrule
\end{tabularx}

\textbf{Protective factors:}

\begin{tabularx}{\textwidth}{p{4cm}p{1.8cm}X}
\toprule
\textbf{Field} & \textbf{Type} & \textbf{Description} \\
\midrule
\trow{protective\_factors}{Dict[str, float]}{social\_support, family\_support, formal\_intervention, psychological\_capital}
\bottomrule
\end{tabularx}

\textbf{Stress tracking:}

\begin{tabularx}{\textwidth}{p{4cm}p{1.8cm}X}
\toprule
\textbf{Field} & \textbf{Type} & \textbf{Description} \\
\midrule
\trow{current\_stress}{float}{Current stress level [0, 1]}
\trow{recent\_stress\_intensity}{float}{Tracks recent stress for PSS-10 response}
\trow{stress\_momentum}{float}{Rate of stress change for predictive updates}
\trow{consecutive\_hindrances}{float}{Hindrance streak length}
\trow{stress\_breach\_count}{int}{Count of stress threshold breaches}
\bottomrule
\end{tabularx}

\textbf{PSS-10 state:}

\begin{tabularx}{\textwidth}{p{4cm}p{1.8cm}X}
\toprule
\textbf{Field} & \textbf{Type} & \textbf{Description} \\
\midrule
\trow{pss10\_responses}{Dict}{Individual PSS-10 item responses}
\trow{stress\_controllability}{float}{Controllability dimension [0, 1]}
\trow{stress\_overload}{float}{Overload dimension [0, 1]}
\trow{pss10}{int}{Total PSS-10 score (0-40)}
\trow{pss10\_smoothed}{float}{Float smoothed value, carried across days}
\trow{stressed}{bool}{Stress classification based on threshold}
\trow{daily\_pss10\_scores}{List[int]}{Scores collected during current day}
\bottomrule
\end{tabularx}

\textbf{Transient keys} (written by one phase, consumed by next):

\begin{tabularx}{\textwidth}{p{4cm}p{1.8cm}XX}
\toprule
\textbf{Field} & \textbf{Type} & \textbf{Producer} & \textbf{Consumer} \\
\midrule
\trowfour{challenge}{float}{Stress Perception}{Resilience Activation}
\trowfour{hindrance}{float}{Stress Perception}{Resilience Activation}
\trowfour{is\_stressed}{bool}{Stress Perception}{Orchestrator branching}
\trowfour{event\_controllability}{float}{Stress Perception}{Resilience Activation}
\trowfour{event\_overload}{float}{Stress Perception}{Resilience Activation}
\bottomrule
\end{tabularx}

Reference:
\href{https://github.com/namakala/abm-simulation/blob/7e40a44f82da76b18910b774cd882c938bc79992/src/python/phases/interfaces.py\#L18-L75}{src/python/phases/interfaces.py:L18-L75}

\section{Two-Loop Orchestration}\label{two-loop-orchestration}

\subsection{Purpose}\label{purpose-2}

Explain how \texttt{Person.step()} orchestrates phases across two
temporal scales within a single simulation day: event-driven responses
and daily homeostatic consolidation.

\begin{itemize}
\tightlist
\item
  \textbf{Frequency:} daily (orchestrator)
\item
  \textbf{Inputs:} AgentState at day start
\item
  \textbf{Outputs:} AgentState at day end
\end{itemize}

\subsection{Algorithm}\label{algorithm}

\begin{verbatim}
FUNCTION step(agent, config, rng):
    // Reset phase outputs
    agent.last_phase_outputs ← {}

    // Build state from agent attributes
    state ← build_state(agent)
    neighbors ← get_neighbor_affects(agent)

    // Subevent loop
    n ← sample_poisson(config.subevents_per_day)
    actions ← random_sequence("stress", "interact", n)
    shuffle(actions)

    challenge_total ← 0
    hindrance_total ← 0
    stress_count ← 0

    FOR EACH action IN actions:
        // Decay support boost
        state.support_boost ← state.support_boost × 0.9

        IF action = "stress":
            // Appraise stress event
            result ← run_stress_perception(state, config, rng)
            state ← apply_delta(state, result.delta)
            challenge_total ← challenge_total + result.delta.challenge
            hindrance_total ← hindrance_total + result.delta.hindrance
            stress_count ← stress_count + 1

            // Activate resilience if stressed
            IF state.is_stressed:
                r2 ← run_resilience_activation(
                    state, config, rng)
                state ← apply_delta(state, r2.delta)
                scores ← state.pss10_scores
                state.pss10_scores ← scores ∪ {state.pss10}

            // Record stress event
            state.stress_events ← state.stress_events ∪ {event}

        ELSE IF action = "interact":
            // Dyadic interaction
            partner ← random_neighbor(agent)
            pstate ← build_state(partner)
            out1, out2 ← interact(state, pstate, config, rng)
            state ← apply_delta(state, out1.delta)
            apply_delta(partner, out2.delta)
            write_back(partner)

            state.interactions ← state.interactions + 1
            IF out1.obs.support_occurred:
                state.support_exchanges ← state.support_exchanges + 1
                sb ← state.support_boost + 0.1
                state.support_boost ← min(1, sb)

    // Normalize daily totals
    IF stress_count > 0:
        challenge_total ← challenge_total ÷ stress_count
        hindrance_total ← hindrance_total ÷ stress_count

    // Daily consolidation
    state ← apply_delta(state,
        affect_dynamics(state, config, rng).delta)
    state ← apply_delta(state,
        resource_allocation(state, config, rng).delta)
    state ← apply_delta(state,
        stress_buffering(state, config, rng).delta)
    state ← apply_delta(state,
        pss10_consolidation(state, config, rng).delta)
    state ← apply_delta(state,
        daily_reset(state, config, rng).delta)

    // Write back
    write_back(agent, state)
\end{verbatim}

\subsection{Parameters}\label{parameters}

\begin{tabularx}{\textwidth}{p{4cm}p{1.8cm}X}
\toprule
\textbf{Name} & \textbf{Description} & \textbf{Default} \\
\midrule
\trow{subevents\_per\_day}{Poisson rate (\(\lambda\)) for daily subevents}{3}
\trow{action\_distribution}{Probability of stress vs interact}{0.5 each}
\trow{support\_boost\_decay}{Per-subevent decay of support boost}{0.9}
\trow{support\_boost\_increment}{Per-support-exchange boost increment}{0.10}
\trow{interaction\_boost\_rate}{Per-interaction resilience boost}{0.005}
\bottomrule
\end{tabularx}

Reference:
\href{https://github.com/namakala/abm-simulation/blob/7e40a44f82da76b18910b774cd882c938bc79992/src/python/agent.py\#L601-L850}{src/python/agent.py:L601-L850}

\section{Phase Implementations}\label{phase-implementations}

\subsection{Stress Perception}\label{stress-perception}

\subsubsection{Purpose}\label{purpose-3}

Transform raw stress events into appraised challenge/hindrance scores
and determine whether the event exceeds the agent's adaptive threshold.
Implements Lazarus primary appraisal.

\begin{itemize}
\tightlist
\item
  \textbf{Frequency:} event\_driven
\item
  \textbf{Inputs:} stress\_controllability, stress\_overload,
  volatility, recent\_stress\_intensity, stress\_momentum
\item
  \textbf{Outputs:} state\_delta: challenge, hindrance, is\_stressed,
  event attrs, updated dimensions
\item
  \textbf{Observation:} event attrs, appraisal values, threshold values
\end{itemize}

\subsubsection{Algorithm}\label{algorithm-1}

\begin{verbatim}
FUNCTION run_stress_perception(state, config, rng):
    // Generate and appraise stress event
    event ← generate_stress_event(rng)
    weights ← {omega_c, omega_o, bias, gamma}
    challenge, hindrance ← apply_weights(event, weights)

    // Compute stress load
    stress_load ← appraised_stress(event, challenge,
        hindrance, config.delta)

    // Evaluate threshold
    threshold ← config.base_threshold
        + config.challenge_scale × challenge
        + config.hindrance_scale × hindrance
    is_stressed ← stress_load ≥ threshold

    // Update stress dimensions
    c, o, inten, mom ← update_dimensions(
        state.controllability, state.overload,
        challenge, hindrance, is_stressed,
        config.volatility, state.stress_intensity,
        state.stress_momentum, state.resilience
    )

    // Build output
    delta ← {
        challenge, hindrance, is_stressed,
        event_controllability: event.controllability,
        event_overload: event.overload,
        stress_controllability: c,
        stress_overload: o,
        stress_intensity: inten,
        stress_momentum: mom
    }
    obs ← {
        event_controllability: event.controllability,
        event_overload: event.overload,
        challenge, hindrance, stress_load,
        threshold, is_stressed
    }
    RETURN {delta, obs}
\end{verbatim}

\subsubsection{Parameters}\label{parameters-1}

\begin{tabularx}{\textwidth}{p{4cm}p{1.8cm}X}
\toprule
\textbf{Name} & \textbf{Description} & \textbf{Default} \\
\midrule
\trow{omega\_c}{Controllability weight in appraisal}{1.0}
\trow{omega\_o}{Overload weight in appraisal}{1.0}
\trow{bias}{Appraisal bias term}{0.0}
\trow{gamma}{Sigmoid steepness in appraisal}{6.0}
\trow{delta}{Polarity effect strength}{0.2}
\trow{base\_threshold}{Base stress threshold}{0.5}
\trow{challenge\_scale}{Challenge threshold scale}{0.15}
\trow{hindrance\_scale}{Hindrance threshold scale}{0.25}
\bottomrule
\end{tabularx}

Reference:
\href{https://github.com/namakala/abm-simulation/blob/7e40a44f82da76b18910b774cd882c938bc79992/src/python/phases/stress_perception.py\#L28-L130}{src/python/phases/stress\_perception.py:L28-L130}

\subsection{Resilience Activation}\label{resilience-activation}

\subsubsection{Purpose}\label{purpose-4}

Determine coping outcome and update resilience after a stressed event.
Implements Lazarus secondary appraisal and coping response. Runs the
full coping pipeline: neighbor influence, coping probability, coping
outcome, affect/resilience/stress changes, PSS-10 generation, resource
cost, and protective factor allocation.

\begin{itemize}
\tightlist
\item
  \textbf{Frequency:} event\_driven
\item
  \textbf{Inputs:} challenge, hindrance (from stress perception),
  affect, resilience, resources, protective\_factors
\item
  \textbf{Outputs:} state\_delta: affect, resilience, current\_stress,
  resources, protective\_factors, pss10, stressed (12 keys)
\item
  \textbf{Observation:} coped\_successfully, coping\_probability,
  resource\_cost
\end{itemize}

\subsubsection{Algorithm}\label{algorithm-2}

\begin{verbatim}
FUNCTION run_resilience_activation(state, config, rng):
    // Determine coping outcome
    a, r, s, coped ← coping_outcome(
        state.affect, state.resilience,
        state.stress, state.challenge,
        state.hindrance, state.neighbors,
        config, rng
    )

    // Update stress dimensions
    c, o, inten, mom ← update_dimensions(
        state.controllability, state.overload,
        state.challenge, state.hindrance,
        coped, config.volatility,
        state.stress_intensity,
        state.stress_momentum, state.resilience
    )

    // Generate PSS-10
    pss10 ← generate_pss10(c, o, inten, mom,
        a, state.resources, r, rng)

    // Resource cost and depletion
    cost ← resource_cost(config.base_cost,
        r, state.challenge, state.hindrance)
    res ← depletion(state.resources, cost, r,
        coped, config)

    // Track consecutive hindrances
    hindrances ← state.consecutive_hindrances
    IF state.hindrance > state.challenge:
        hindrances ← hindrances + 1
    ELSE:
        hindrances ← 0

    // PF allocation if coped
    pf ← copy(state.protective_factors)
    IF coped:
        res ← clamp(res + 0.03, 0, 1)
        alloc ← allocate_pf(
            res × config.pf_alloc_fraction,
            r, state.baseline_resilience, pf, rng
        )
        pf ← update_pf(pf, alloc, r)
        res ← clamp(res - sum(alloc), 0, 1)

    // Build output
    delta ← {
        affect: a, resilience: r,
        stress: s,
        controllability: c, overload: o,
        resources: res, protective_factors: pf,
        consecutive_hindrances: hindrances,
        stress_breach_count: state.breach_count + 1,
        pss10: pss10.score,
        pss10_responses: pss10.responses,
        stressed: pss10.stressed
    }
    obs ← {
        coped, coping_probability,
        resilience_effect, resource_cost: cost,
        delta_stress: s - state.stress,
        delta_affect: a - state.affect
    }
    RETURN {delta, obs}
\end{verbatim}

\subsubsection{Parameters}\label{parameters-2}

\begin{tabularx}{\textwidth}{p{4cm}p{1.8cm}X}
\toprule
\textbf{Name} & \textbf{Description} & \textbf{Default} \\
\midrule
\trow{base\_resource\_cost}{Cost per coping attempt}{0.1}
\trow{pf\_allocation\_fraction}{Resources allocated to PF on success}{0.15}
\trow{resource\_reward}{Flat reward for successful coping}{0.03}
\bottomrule
\end{tabularx}

Reference:
\href{https://github.com/namakala/abm-simulation/blob/7e40a44f82da76b18910b774cd882c938bc79992/src/python/phases/resilience_activation.py\#L42-L252}{src/python/phases/resilience\_activation.py:L42-L252}

\subsection{Social Interaction}\label{social-interaction}

\subsubsection{Purpose}\label{purpose-5}

Dyadic interaction between two agents. Converges affect and resilience
with negativity bias (negative influence 1.5x stronger), detects support
from convergence magnitude, and applies resource exchange state machine
based on mutual stress states.

\begin{itemize}
\tightlist
\item
  \textbf{Frequency:} event\_driven
\item
  \textbf{Inputs:} self\_state, partner\_state (full AgentState)
\item
  \textbf{Outputs:} Tuple{[}PhaseOutput, PhaseOutput{]}: delta values
  (not absolute)
\item
  \textbf{Observation:} support\_occurred
\end{itemize}

\subsubsection{Algorithm}\label{algorithm-3}

\begin{verbatim}
FUNCTION interact(self_state, partner_state, config, rng):
    // Affect convergence with negativity bias
    dA_self ← config.rate × partner.affect
    dA_partner ← config.rate × self.affect
    IF dA_self < 0:
        dA_self ← dA_self × 1.5
    IF dA_partner < 0:
        dA_partner ← dA_partner × 1.5

    // Resilience convergence
    dR_self ← config.resilience_rate × partner.affect
    dR_partner ← config.resilience_rate × self.affect

    // Support detection
    total ← |dA_self| + |dA_partner|
        + |dR_self| + |dR_partner|
    support ← total > config.threshold

    // Resource exchange
    IF self.stressed AND partner.stressed:
        IF support: boost both
        ELSE: cost both
    ELSE IF self.stressed:
        IF support: boost self
        ELSE: cost self
    ELSE IF partner.stressed:
        IF support: boost partner
        ELSE: cost partner
    ELSE:
        small_boost(self, partner)

    RETURN (self_out, partner_out)
\end{verbatim}

\subsubsection{Parameters}\label{parameters-3}

\begin{tabularx}{\textwidth}{p{4cm}p{1.8cm}X}
\toprule
\textbf{Name} & \textbf{Description} & \textbf{Default} \\
\midrule
\trow{influence\_rate}{Affect convergence rate}{0.05}
\trow{resilience\_influence}{Resilience convergence rate}{0.05}
\trow{negativity\_bias}{Negative influence multiplier}{1.5}
\trow{support\_threshold}{Convergence threshold for support detection}{0.1}
\trow{boost}{Resource boost on support exchange}{0.05}
\trow{cost}{Resource cost without support}{0.03}
\bottomrule
\end{tabularx}

Reference:
\href{https://github.com/namakala/abm-simulation/blob/7e40a44f82da76b18910b774cd882c938bc79992/src/python/phases/interaction.py\#L48-L190}{src/python/phases/interaction.py:L48-L190}

\subsection{Resource Allocation}\label{resource-allocation}

\subsubsection{Purpose}\label{purpose-6}

Regenerates resources and allocates to protective factors via softmax.
Resource regeneration is linear in the deficit, modulated by affect and
resilience. Allocation uses softmax with temperature for
bounded-rational distribution across social support, family support,
formal intervention, and psychological capital.

\begin{itemize}
\tightlist
\item
  \textbf{Frequency:} daily
\item
  \textbf{Inputs:} resources, affect, resilience, protective\_factors
\item
  \textbf{Outputs:} state\_delta: resources, protective\_factors
\item
  \textbf{Observation:} regeneration\_amount, allocation\_weights
\end{itemize}

\subsubsection{Algorithm}\label{algorithm-4}

\begin{verbatim}
FUNCTION run_resource_allocation(state, config, rng):
    // Resource regeneration
    a_mult ← 1 + 0.5 × max(0, state.affect)
    r_mult ← 1 + 0.3 × state.resilience
    regen ← config.base_rate × (1 - state.resources)
        × a_mult × r_mult

    // Softmax allocation
    available ← state.resources + regen
    spendable ← available × (1 - config.preserve_frac)
    weights ← softmax(state.efficacies / config.temp)
    alloc ← spendable × weights

    // Update efficacies (diminishing returns)
    FOR each factor f:
        de ← alloc[f] × config.rate
            × (1 - state.efficacies[f])
        state.efficacies[f] ← min(1,
            state.efficacies[f] + de)

    // Remaining resources
    new_res ← config.preserve_frac × available
        + (spendable - total(alloc))

    RETURN {delta: {resources: new_res,
        protective_factors: state.efficacies},
        obs: {regen, alloc}}
\end{verbatim}

\subsubsection{Parameters}\label{parameters-4}

\begin{tabularx}{\textwidth}{p{4cm}p{1.8cm}X}
\toprule
\textbf{Name} & \textbf{Description} & \textbf{Default} \\
\midrule
\trow{base\_regeneration}{Daily regeneration rate}{0.5}
\trow{preservable\_fraction}{Fraction of resources preserved}{0.2}
\trow{softmax\_temperature}{Allocation temperature}{1.0}
\trow{protective\_improvement\_rate}{PF efficacy update rate}{0.5}
\bottomrule
\end{tabularx}

Reference:
\href{https://github.com/namakala/abm-simulation/blob/7e40a44f82da76b18910b774cd882c938bc79992/src/python/phases/resource_allocation.py\#L141-L219}{src/python/phases/resource\_allocation.py:L141-L219}

\subsection{Stress Buffering}\label{stress-buffering}

\subsubsection{Purpose}\label{purpose-7}

Applies protective-factor resilience boost and resource mediation of
stress buffering using Baron \& Kenny mediation framework: a-path
(stress → resources), b-path (resources → buffering), c'-path (stress →
buffering \textbar{} resources).

\begin{itemize}
\tightlist
\item
  \textbf{Frequency:} daily
\item
  \textbf{Inputs:} resilience, baseline\_resilience,
  protective\_factors, current\_stress, resources
\item
  \textbf{Outputs:} state\_delta: resilience, resources
\item
  \textbf{Observation:} pf\_boost, buffering\_strength, indirect\_effect
\end{itemize}

\subsubsection{Algorithm}\label{algorithm-5}

\begin{verbatim}
FUNCTION run_stress_buffering(state, config, rng):
    // PF boost to resilience
    need ← max(0, state.baseline - state.resilience)
    boost ← Σ(efficacy × need × config.rate)
    boost ← min(boost, need)
    new_r ← state.resilience + boost

    // Resource mediation (a-path)
    eff_stress ← state.stress × (1 + 0.2 × state.overload)
    depletion ← config.a_coeff × eff_stress
    new_res ← state.resources + depletion

    // Buffering strength (b-path + c'-path)
    buf ← max(0,
        config.b_coeff × new_res
        + config.c_coeff × state.stress)

    RETURN {delta: {resilience: new_r,
        resources: new_res},
        obs: {boost, buf}}
\end{verbatim}

\subsubsection{Parameters}\label{parameters-5}

\begin{tabularx}{\textwidth}{p{4cm}p{1.8cm}X}
\toprule
\textbf{Name} & \textbf{Description} & \textbf{Default} \\
\midrule
\trow{boost\_rate}{PF resilience boost rate}{0.1}
\trow{a\_coefficient}{Stress→resources path coefficient}{0.3}
\trow{b\_coefficient}{Resources→buffering path coefficient}{0.7}
\trow{c\_prime\_coefficient}{Stress→buffering path coefficient}{1.0}
\bottomrule
\end{tabularx}

Reference:
\href{https://github.com/namakala/abm-simulation/blob/7e40a44f82da76b18910b774cd882c938bc79992/src/python/phases/stress_buffering.py\#L42-L148}{src/python/phases/stress\_buffering.py:L42-L148}

\subsection{Affect Dynamics}\label{affect-dynamics}

\subsubsection{Purpose}\label{purpose-8}

Applies daily affect and resilience dynamics including peer influence,
event appraisal effects, homeostasis, protective factor boosts, and
social resilience optimisation. Runs once per day after the subevent
loop.

\begin{itemize}
\tightlist
\item
  \textbf{Frequency:} daily
\item
  \textbf{Inputs:} affect, baseline\_affect, resilience, resources,
  current\_stress, neighbor\_affects
\item
  \textbf{Outputs:} state\_delta: affect, resilience,
  consecutive\_hindrances
\item
  \textbf{Observation:} neighbor\_affects\_summary, protective\_boost
\end{itemize}

\subsubsection{Algorithm}\label{algorithm-6}

\begin{verbatim}
FUNCTION run_affect_dynamics(state, config, rng):
    // Affect: homeostasis + peers + appraisal
    new_a ← update_affect(state.affect,
        state.baseline_affect, state.neighbors,
        state.challenge, state.hindrance,
        state.stress, state.resources)

    // Resilience dynamics + PF boost
    new_r ← update_resilience(state.resilience,
        state.consecutive_hindrances)
    pf_boost ← get_pf_boost(state.protective_factors,
        state.baseline, new_r)
    new_r ← min(1, new_r + pf_boost)

    // Social resilience optimisation
    new_r ← social_resilience(new_r,
        state.interactions, state.support_exchanges,
        state.resources, state.baseline_resilience,
        state.protective_factors, rng)

    // Interaction-frequency boost
    new_r ← new_r + state.interactions × 0.005

    // Hindrance decay
    new_h ← max(0,
        state.consecutive_hindrances - config.decay)

    // Homeostatic adjustment
    new_a ← homeostatic(state.baseline_affect,
        new_a, config.affect_rate)
    new_r ← homeostatic(state.baseline_resilience,
        new_r, config.resilience_rate)

    RETURN {delta: {affect: new_a,
        resilience: new_r,
        consecutive_hindrances: new_h},
        obs: {}}
\end{verbatim}

\subsubsection{Parameters}\label{parameters-6}

\begin{tabularx}{\textwidth}{p{4cm}p{1.8cm}X}
\toprule
\textbf{Name} & \textbf{Description} & \textbf{Default} \\
\midrule
\trow{affect\_homeostatic\_rate}{Affect return rate}{0.5}
\trow{resilience\_homeostatic\_rate}{Resilience return rate}{0.5}
\trow{interaction\_boost\_rate}{Per-interaction resilience boost}{0.005}
\trow{hindrance\_decay\_rate}{Daily hindrance decay}{0.05}
\bottomrule
\end{tabularx}

Reference:
\href{https://github.com/namakala/abm-simulation/blob/7e40a44f82da76b18910b774cd882c938bc79992/src/python/agent.py\#L67-L210}{src/python/agent.py:L67-L210}

\subsection{PSS-10 Consolidation}\label{pss-10-consolidation}

\subsubsection{Purpose}\label{purpose-9}

Consolidates daily PSS-10 scores, updates stress level via exponential
smoothing, and applies post-hoc adjustments for resilience coupling,
resource coupling, and mood-congruent appraisal.

\begin{itemize}
\tightlist
\item
  \textbf{Frequency:} daily
\item
  \textbf{Inputs:} daily\_pss10\_scores, current\_stress,
  stress\_controllability, stress\_overload
\item
  \textbf{Outputs:} state\_delta: pss10, pss10\_smoothed,
  current\_stress, stressed, daily\_pss10\_scores (cleared)
\item
  \textbf{Observation:} avg\_pss10, num\_events
\end{itemize}

\subsubsection{Algorithm}\label{algorithm-7}

\begin{verbatim}
FUNCTION run_pss10_consolidation(state, config, rng):
    // Update dimensions from PSS-10 feedback
    c, o ← update_dimensions_from_pss10(
        state.controllability, state.overload,
        state.pss10_responses, state.resources)

    // Compute stress from dimensions
    new_s ← stress_from_dimensions(c, o,
        state.affect, state.resources,
        state.resilience)
    s ← 0.5 × new_s + 0.5 × state.stress

    // Consolidate daily scores
    IF state.daily_scores ≠ []:
        consolidated ← mean(state.daily_scores)
    ELSE:
        consolidated ← regenerate(c, o)

    // Exponential smoothing
    α ← config.smoothing_alpha
    smoothed ← α × consolidated + (1 - α) × state.pss10_smoothed

    // Post-hoc adjustments
    final ← smoothed + config.bias
        - 0.5 × config.coupling
        × (state.resilience - 0.5)
    final ← 20 + (final - 20) × 1.2
    final ← final + (0.5 - state.resources) × 12
    stressed ← final ≥ config.threshold

    // Mood-congruent appraisal
    adj ← -state.affect × config.affect_adj

    RETURN {delta: {pss10: final + adj,
        pss10_smoothed: smoothed,
        stress: s, controllability: c,
        overload: o, stressed,
        daily_scores: []},
        obs: {}}
\end{verbatim}

\subsubsection{Parameters}\label{parameters-7}

\begin{tabularx}{\textwidth}{p{4cm}p{1.8cm}X}
\toprule
\textbf{Name} & \textbf{Description} & \textbf{Default} \\
\midrule
\trow{pss10\_smoothing\_alpha}{Exponential smoothing factor}{0.3}
\trow{pss10\_resilience\_coupling}{Resilience penalty coefficient}{3.5}
\trow{pss10\_threshold}{Clinical cutoff score}{27}
\trow{pss10\_affect\_adjustment}{Mood-congruent appraisal factor}{0.5}
\bottomrule
\end{tabularx}

Reference:
\href{https://github.com/namakala/abm-simulation/blob/7e40a44f82da76b18910b774cd882c938bc79992/src/python/agent.py\#L212-L348}{src/python/agent.py:L212-L348}

\subsection{Daily Reset}\label{daily-reset}

\subsubsection{Purpose}\label{purpose-10}

Clears daily tracking counters, applies affect reset toward baseline,
and decays stress and consecutive hindrances. Prepares agent state for
the next day's events.

\begin{itemize}
\tightlist
\item
  \textbf{Frequency:} daily
\item
  \textbf{Inputs:} affect, baseline\_affect, current\_stress,
  consecutive\_hindrances, daily counters
\item
  \textbf{Outputs:} state\_delta: affect, stress, counters cleared
\item
  \textbf{Observation:} stress\_summary
\end{itemize}

\subsubsection{Algorithm}\label{algorithm-8}

\begin{verbatim}
FUNCTION run_daily_reset(state, config, rng):
    // Affect reset toward baseline
    rate ← scale_rate(config.affect_rate,
        state.resources, state.stress)
    new_a ← affect_reset(state.affect,
        state.baseline_affect, rate)

    // Stress decay
    new_s ← stress_decay(state.stress, rate)

    // Hindrance decay
    new_h ← max(0,
        state.consecutive_hindrances - 0.05)

    // Stress summary
    events ← state.stress_events
    summary ← {
        avg: mean(events.stress_level),
        max: max(events.stress_level),
        n: len(events),
        coped: mean(events.coped)
    }

    // Clear daily counters
    state.interactions ← 0
    state.support_exchanges ← 0
    state.stress_events ← []
    state.pss10_scores ← []

    RETURN {delta: {affect: new_a,
        stress: new_s,
        consecutive_hindrances: new_h,
        last_reset_day: state.day},
        obs: {summary}}
\end{verbatim}

\subsubsection{Parameters}\label{parameters-8}

\begin{tabularx}{\textwidth}{p{4cm}p{1.8cm}X}
\toprule
\textbf{Name} & \textbf{Description} & \textbf{Default} \\
\midrule
\trow{affect\_homeostatic\_rate}{Affect return rate}{0.5}
\trow{hindrance\_decay\_rate}{Daily hindrance decay}{0.05}
\bottomrule
\end{tabularx}

Reference:
\href{https://github.com/namakala/abm-simulation/blob/7e40a44f82da76b18910b774cd882c938bc79992/src/python/agent.py\#L350-L412}{src/python/agent.py:L350-L412}

\section{Rationale for Cumulative Block
Experiment}\label{rationale-for-cumulative-block-experiment}

\subsection{Purpose}\label{purpose-11}

Explain why stages 1-7 selectively enable phases and how this
demonstrates the modularity claim made in the main text.

\begin{itemize}
\tightlist
\item
  \textbf{Frequency:} N/A (analytical design)
\item
  \textbf{Inputs:} N/A
\item
  \textbf{Outputs:} N/A
\end{itemize}

\subsection{Design}\label{design}

The cumulative block experiment assembles the model block by block
across seven stages, each running the same seeded population (N = 100
agents, 90 days). Each stage adds one building block on top of the
previous ones:

\begin{tabularx}{\textwidth}{p{1.5cm}p{4cm}X}
\toprule
\textbf{Stage} & \textbf{Blocks Active} & \textbf{Purpose} \\
\midrule
1 & Initialization only & Baseline population state \\
2 & + Stress perception & Isolate appraisal effects \\
3 & + Resilience activation & Isolate coping effects \\
4 & + Interaction & Isolate social contagion \\
5 & + Resource allocation & Isolate resource regeneration \\
6 & + Stress buffering & Isolate protective buffering \\
7 & + Full model & Complete homeostatic cycle \\
\bottomrule
\end{tabularx}

\subsection{Interpretation}\label{interpretation}

Intermediate stages (2-6) are deliberately partial: they enable
stress-adding mechanisms without the restorative blocks of the daily
cycle. This isolation demonstrates that:

\begin{enumerate}
\def\labelenumi{\arabic{enumi}.}
\tightlist
\item
  \textbf{Stress perception alone} (stage 2) updates appraisal
  dimensions without affecting outcome states, so stages 1 and 2
  coincide
\item
  \textbf{Resilience activation} (stage 3) introduces coping, causing
  population depletion without regeneration
\item
  \textbf{Interaction} (stage 4) propagates affect/resilience through
  the network with negativity bias, dragging both toward lower bounds
\item
  \textbf{Resource allocation} (stage 5) restores the resource pool but
  cannot move affect or resilience
\item
  \textbf{Stress buffering} (stage 6) begins pushing resilience back
  toward baseline
\end{enumerate}

Stage 7, the complete model, activates the remaining blocks: affect
dynamics, PSS-10 consolidation, daily reset, and network adaptation.
This completes the two-loop daily cycle. The homeostatic pull toward
baseline affect and resilience, the daily stress decay, and adaptive
rewiring toward more similar peers jointly restore equilibrium.

\subsection{Analytical vs Runtime
Stacking}\label{analytical-vs-runtime-stacking}

The cumulative block experiment demonstrates \textbf{analytical
stacking}: selectively enabling/disabling phases to isolate
contributions. At runtime, the simulation uses \textbf{sequential
composition}: all phases execute every day in a pipeline, with each
phase consuming the previous phase's output.

This distinction is important: the modularity claim is validated by the
cumulative experiment, while the runtime behavior uses the full two-loop
orchestration described in S3.

See: article.qmd Results → ``Mechanism contribution: cumulative
building-block build-up''

\section{Network Topology}\label{network-topology}

\subsection{Purpose}\label{purpose-12}

Social network uses Watts-Strogatz small-world topology with high
clustering and short characteristic path lengths. Network adaptation
rewires edges based on stress breach counts and homophily similarity.

\begin{itemize}
\tightlist
\item
  \textbf{Frequency:} model-level (network adaptation runs daily)
\item
  \textbf{Inputs:} agent stress\_breach\_count, affect, resilience
\item
  \textbf{Outputs:} modified graph edges
\end{itemize}

\subsection{Algorithm}\label{algorithm-9}

\begin{verbatim}
FUNCTION build_network(N, k, p, rng):
    k ← min(k, N - 1)
    RETURN watts_strogatz(N, k, p, rng)

FUNCTION adapt_network(G, agents, config, rng):
    FOR each node WHERE breach_count ≥ threshold:
        worst ← most dissimilar neighbor
        prob ← sigmoid(similarity, support, homophily)
        IF rng > prob:
            best ← most similar non-neighbor
            rewire(node, worst, best)
    RETURN G
\end{verbatim}

\subsection{Parameters}\label{parameters-9}

\begin{tabularx}{\textwidth}{p{4cm}p{1.8cm}X}
\toprule
\textbf{Name} & \textbf{Description} & \textbf{Default} \\
\midrule
\trow{N}{Number of agents}{20}
\trow{k}{Mean degree (must be even, < N)}{4}
\trow{p}{Rewiring probability}{0.1}
\trow{threshold}{Breach count for adaptation}{3}
\trow{homophily}{Similarity weight}{0.7}
\bottomrule
\end{tabularx}

Reference:
\href{https://github.com/namakala/abm-simulation/blob/7e40a44f82da76b18910b774cd882c938bc79992/src/python/network_utils.py\#L16-L245}{src/python/network\_utils.py:L16-L245}

\section{Data Collection}\label{data-collection}

\subsection{Purpose}\label{purpose-13}

Mesa's DataCollector captures model-level and agent-level metrics using
named reporter functions. Model reporters track population averages;
agent reporters track individual trajectories.

\begin{itemize}
\tightlist
\item
  \textbf{Frequency:} model-level (collected each step)
\item
  \textbf{Inputs:} agent attributes, model attributes
\item
  \textbf{Outputs:} CSV export
\end{itemize}

\subsection{Agent Reporters}\label{agent-reporters}

\begin{verbatim}
AGENT_REPORTERS = {
    pss10, resilience, affect, resources, current_stress,
    stress_controllability, stress_overload, consecutive_hindrances,
    coping_success, challenge_appraisal, hindrance_appraisal,
    interaction_frequency, stressed, support_boost
}
\end{verbatim}

\subsection{Model Reporters}\label{model-reporters}

\begin{verbatim}
MODEL_REPORTERS = {
    avg_pss10, avg_resilience, avg_affect, coping_success_rate,
    avg_resources, avg_stress, social_support_rate, network_density,
    stress_prevalence, low_resilience, high_resilience,
    avg_challenge, avg_hindrance, challenge_hindrance_ratio
}
\end{verbatim}

\subsection{Phase-Level
Instrumentation}\label{phase-level-instrumentation}

The \texttt{\_last\_phase\_outputs} dictionary stores the
\texttt{PhaseOutput} from each phase executed during a step. This
enables phase-specific metric extraction for the cumulative block
experiment and debugging.

\begin{verbatim}
last_phase_outputs = {
    stress_perception: PhaseOutput,
    resilience_activation: PhaseOutput,
    interaction_self: PhaseOutput,
    interaction_partner: PhaseOutput,
    affect_dynamics: PhaseOutput,
    resource_allocation: PhaseOutput,
    stress_buffering: PhaseOutput,
    pss10_consolidation: PhaseOutput,
    daily_reset: PhaseOutput,
}
\end{verbatim}

Reference:
\href{https://github.com/namakala/abm-simulation/blob/7e40a44f82da76b18910b774cd882c938bc79992/src/python/reporters.py\#L289-L333}{src/python/reporters.py:L289-L333}

Reference:
\href{https://github.com/namakala/abm-simulation/blob/7e40a44f82da76b18910b774cd882c938bc79992/src/python/model.py\#L114-L133}{src/python/model.py:L114-L133}




\end{document}
